% paper_gravity_control_sigma.tex
% Gravitational Sigma Engineering: Feasibility of Local Spacetime
% Modification via Quantum Information Control
% Author: Sheng-Kai Huang
% Compile: pdflatex paper_gravity_control_sigma.tex
% Target: ~5 pages, PRD style

\documentclass[prd,twocolumn,superscriptaddress,longbibliography,floatfix]{revtex4-2}

\usepackage{amsmath}
\usepackage{amssymb}
\usepackage{amsfonts}
\usepackage{amsthm}
\usepackage{bm}
\usepackage{hyperref}
\usepackage{booktabs}
\usepackage{graphicx}

\newtheorem{theorem}{Theorem}
\newtheorem{corollary}[theorem]{Corollary}
\newtheorem*{proposition}{Proposition}

% Shared macros (consistent with Papers 1--6)
\newcommand{\kB}{k_{\mathrm{B}}}
\newcommand{\Tr}{\mathrm{Tr}}
\newcommand{\Petz}{\mathcal{R}}
\newcommand{\chan}{\mathcal{N}}
\newcommand{\ket}[1]{|#1\rangle}
\newcommand{\bra}[1]{\langle#1|}
\newcommand{\tauasym}{\tau}
\newcommand{\Sgrav}{\Sigma_{\mathrm{grav}}}
\newcommand{\rs}{r_{\mathrm{s}}}

% Paper-specific macros
\newcommand{\DSigma}{\Delta\Sigma}
\newcommand{\rhoK}{\rho_{\mathrm{K}}}
\newcommand{\rhoCas}{\rho_{\mathrm{Cas}}}
\newcommand{\vF}{v_{\mathrm{F}}}
\newcommand{\trel}{\tau_{\mathrm{relax}}}

% Epistemic tags
\newcommand{\known}{\textsc{[known]}}
\newcommand{\derived}{\textsc{[derived]}}
\newcommand{\assumed}{\textsc{[assumed]}}
\newcommand{\new}{\textsc{[new]}}
\newcommand{\specmark}{\textsc{[speculative]}}

\begin{document}

\title{Gravitational Sigma Engineering: Feasibility of Local Spacetime
Modification via Quantum Information Control}

\author{Sheng-Kai Huang}
\email{akai@fawstudio.com}
\affiliation{Independent Researcher}

\date{\today}

\begin{abstract}
Within the $\Sigma$--$\tauasym$ framework, the gravitational entropy
production $\Sgrav = -\ln(-g_{00})$ is identified as a fundamental
field governing the gravitational channel.  If $\Sgrav$ can be
modified locally \emph{without} adding mass, one obtains ``gravity
control'' in the strict sense of altering the effective metric.  We
evaluate four candidate mechanisms quantitatively: (i)~Khronon field
condensate displacement ($\DSigma \sim 10^{-97}$), (ii)~Casimir
vacuum energy ($\DSigma \sim 10^{-46}$), (iii)~superconductor
condensation energy ($\DSigma \sim 10^{-34}$), and (iv)~nanoparticle
matter-wave interferometry ($\DSigma \sim 10^{-10}$ for
next-generation experiments).  All estimates are compared against
experimental detection thresholds: torsion-balance sensitivity
$\DSigma \sim 10^{-15}$, atom interferometry $\DSigma \sim 10^{-20}$,
and LIGO-class detectors $\DSigma \sim 10^{-22}$.  We find that
approaches~(i)--(iii) are quantitatively excluded for laboratory-scale
gravity modification, while approach~(iv) offers a realistic path to
\emph{detecting} $\Sgrav$ at the quantum scale within the coming
decade.  We identify the quantum-coherence channel of superconductors
as an open theoretical question deserving further investigation.
\end{abstract}

\maketitle

%======================================================================
\section{Introduction}
\label{sec:intro}
%======================================================================

General relativity describes gravity through the spacetime metric
$g_{\mu\nu}$, whose curvature encodes the gravitational field
\known.  In the companion papers of this series, we established a
quantum-information interpretation: the temporal asymmetry parameter
$\tauasym = 1 - F$, where $F$ is the Petz recovery fidelity, quantifies
the irreversibility of the gravitational channel
\cite{Huang2026paper1}.  The Petz recovery bound
\begin{equation}
  F \geq e^{-\Sigma/2}
  \label{eq:petz-bound}
\end{equation}
relates the recovery fidelity to the entropy production $\Sigma$
\known~\cite{Petz1986,Petz1988,Junge2018,Faulkner2022}.
In Paper~2 \cite{Huang2026paper2}, we proved the Channel Theorem:
for a static gravitational background, the extensive entropy
production of the gravitational channel is
\begin{equation}
  \Sgrav = -\ln(-g_{00}) \,,
  \label{eq:sigma-g00}
\end{equation}
which for the Schwarzschild exterior gives
$\Sgrav = -\ln(1 - \rs/r)$ \derived.

A key consequence is the Fisher information principle: demanding that
$\Sgrav$ be a harmonic function in vacuum,
$\nabla^2 \Sgrav = 0$, yields the exponential metric
\begin{equation}
  g_{00} = -e^{-\rs/r}\,, \qquad g_{rr} = e^{+\rs/r}\,,
  \label{eq:exp-metric}
\end{equation}
which is horizon-free and agrees with all current weak-field tests
\derived~\cite{Huang2026paper2}.

This raises a provocative question: if $\Sgrav$ is a fundamental
field rather than merely a derived quantity, can it be modified
\emph{locally} without placing mass at that location?  Such a
modification would constitute ``gravity control'' in the literal
sense---altering the effective metric $g_{00}$ at a chosen point.
In this paper, we evaluate this possibility quantitatively, computing
$\DSigma$ achievable by four distinct physical mechanisms and
comparing each against experimental detection thresholds.  We
emphasize that honest assessment requires showing both what works
and what does not.

%======================================================================
\section{Defining Gravity Control}
\label{sec:definition}
%======================================================================

We define \emph{gravity control} operationally:

\begin{proposition}[Gravity control]
A change $\DSigma \neq 0$ at spatial location $\bm{r}$, produced
by a mechanism that does not introduce additional stress-energy
$\Delta T_{\mu\nu}(\bm{r}) = 0$ at that location.
\end{proposition}

\noindent From Eq.~\eqref{eq:sigma-g00}, such a change implies a
local metric modification
\begin{equation}
  g_{00} \;\to\; g_{00}\,e^{-\DSigma}\,,
  \label{eq:metric-shift}
\end{equation}
and a corresponding change in the local gravitational acceleration
\begin{equation}
  \Delta a = \frac{c^2}{2}\,\frac{\partial(\DSigma)}{\partial r}
  \label{eq:accel-shift}
\end{equation}
for a spatially varying $\DSigma(r)$ \derived.

\subsection{Detection thresholds}

The sensitivity of current and planned experiments sets the minimum
detectable $\DSigma$ \known:

\begin{itemize}
  \item \textbf{Torsion balance} (E\"ot-Wash
    group~\cite{Adelberger2009}): acceleration sensitivity
    $\delta a \sim 10^{-15}\,\mathrm{m/s^2}$.  Over a baseline
    $\ell \sim 1\,\mathrm{cm}$, this corresponds to
    \begin{equation}
      \DSigma_{\min}^{\mathrm{torsion}} \sim
      \frac{2\,\delta a\,\ell}{c^2}
      = \frac{2 \times 10^{-15} \times 10^{-2}}{9 \times 10^{16}}
      \approx 2 \times 10^{-34}\,.
      \label{eq:torsion-threshold}
    \end{equation}
    However, for a \emph{uniform} $\DSigma$ over a macroscopic region
    $L \sim 1\,\mathrm{m}$, the directly measurable quantity is
    the redshift $\Delta\nu/\nu = \DSigma/2$, giving a threshold
    $\DSigma \sim 10^{-15}$ from optical clock comparisons
    \cite{Bothwell2022}.
  \item \textbf{Atom interferometer}~\cite{Asenbaum2020}:
    phase sensitivity $\Delta\phi \sim 10^{-4}\,\mathrm{rad}$ with
    interrogation time $T \sim 1\,\mathrm{s}$ gives
    $\delta a \sim 10^{-12}\,\mathrm{m/s^2}$, but over long baselines
    $L \sim 10\,\mathrm{m}$ the effective $\DSigma$ threshold reaches
    $\sim 10^{-20}$.
  \item \textbf{LIGO}~\cite{LIGOScientific2015}: strain sensitivity
    $h \sim 10^{-22}$ corresponds to transient $\DSigma \sim 10^{-22}$,
    but only for oscillatory signals in the detector bandwidth.
\end{itemize}

%======================================================================
\section{Approach 1: Khronon Field Condensate}
\label{sec:khronon}
%======================================================================

In Paper~3 \cite{Huang2026paper3}, the Khronon action with ghost
condensation \known~\cite{ArkaniHamed2004} is
\begin{equation}
  K(Q) = \mu^2 (Q - 1)^2\,,
  \label{eq:KQ}
\end{equation}
where $Q = 1/N_\phi$ is the inverse lapse of the Khronon field and
$\mu_0 = H_0/c$ is the single free parameter \derived.  The
condensate sits at $Q_0 = 1$; a displacement $Q = Q_0 + \delta Q$
produces \derived
\begin{equation}
  \DSigma = 2\ln(1 + \delta Q) \approx 2\,\delta Q
  \quad (\delta Q \ll 1)\,.
  \label{eq:dsigma-khronon}
\end{equation}
The energy density required to sustain this displacement is
\begin{equation}
  \rhoK = \frac{\mu^2 c^2}{8\pi G}\,\delta Q\,(2 + \delta Q)\,.
  \label{eq:rho-khronon}
\end{equation}
With $\mu = H_0/c \approx 7.3 \times 10^{-27}\,\mathrm{m}^{-1}$
and $G = 6.674 \times 10^{-11}\,\mathrm{m^3\,kg^{-1}\,s^{-2}}$:
\begin{align}
  \mu^2 c^2 &= (7.3 \times 10^{-27})^2 \times (3 \times 10^8)^2
  \nonumber \\
  &= 5.3 \times 10^{-53} \times 9 \times 10^{16}
  \nonumber \\
  &= 4.8 \times 10^{-36}\;\mathrm{kg\,m^{-1}\,s^{-2}}\,.
  \label{eq:mu2c2}
\end{align}
For a target $\DSigma = 10^{-15}$ (optical-clock detectable),
we need $\delta Q = 5 \times 10^{-16}$, giving
\begin{equation}
  \rhoK = \frac{4.8 \times 10^{-36} \times 10^{-15}}
  {8\pi \times 6.674 \times 10^{-11}}
  \approx \frac{4.8 \times 10^{-51}}{1.67 \times 10^{-9}}
  \approx 2.9 \times 10^{-42}\;\frac{\mathrm{kg}}{\mathrm{m}^3}\,.
  \label{eq:rho-khronon-num}
\end{equation}
This is not the issue---the energy is negligible.  The problem is
the \emph{inverse}: the Khronon condensate at cosmological coupling
$\mu = H_0/c$ only produces $\DSigma \sim 10^{-15}$ at the cost
of displacing $Q$ on cosmological scales.  In the laboratory,
with no cosmological horizon to set the boundary condition,
$\delta Q$ relaxes to zero exponentially on the scale
$\lambda_K = 1/\mu \sim c/H_0 \sim 10^{26}\,\mathrm{m}$ (the
Hubble radius) \derived.

The achievable $\DSigma$ in a lab of size $L \sim 1\,\mathrm{m}$
is suppressed by $(L/\lambda_K)^2$:
\begin{equation}
  \DSigma_{\mathrm{lab}} \sim \DSigma_{\mathrm{cosmo}}
  \times \left(\frac{L}{\lambda_K}\right)^2
  \sim 1 \times \left(\frac{1}{10^{26}}\right)^2
  = 10^{-52}\,.
  \label{eq:khronon-lab}
\end{equation}
Comparing with the Khronon stress-energy contribution over the
lab volume $V = L^3$:
\begin{equation}
  \DSigma \sim \frac{8\pi G\,\rhoK\,L^2}{c^4}
  = \frac{8\pi \times 6.67 \times 10^{-11} \times \rhoK \times 1}
  {8.1 \times 10^{33}}\,.
  \label{eq:dsigma-energy}
\end{equation}
To reach $\DSigma = 10^{-15}$, one would need
$\rhoK \sim 10^{19}\,\mathrm{kg/m^3}$---nuclear density.  But the
Khronon condensate at $\mu = H_0/c$ produces
$\rhoK \sim 10^{-26}\,\mathrm{kg/m^3}$ (the critical density),
not nuclear density.

\textbf{Conclusion:} Direct Khronon condensation at the cosmological
coupling is quantitatively excluded for laboratory gravity
modification.  Any locally engineered $\DSigma$ via Khronon
displacement would require either (a)~a coupling $\mu \gg H_0/c$
(violating the cosmological identification) or (b)~boundary conditions
mimicking a cosmological horizon at lab scales (physically unrealizable).

%======================================================================
\section{Approach 2: Casimir Vacuum Energy}
\label{sec:casimir}
%======================================================================

The Casimir effect \known~\cite{Casimir1948,Lamoreaux1997} produces a
negative energy density between parallel conducting plates separated
by distance $d$:
\begin{equation}
  \rhoCas = -\frac{\pi^2\hbar c}{720\,d^4}\,.
  \label{eq:rho-casimir}
\end{equation}
For $d = 100\,\mathrm{nm} = 10^{-7}\,\mathrm{m}$:
\begin{align}
  \rhoCas &= -\frac{\pi^2 \times 1.055 \times 10^{-34}
  \times 3 \times 10^8}{720 \times (10^{-7})^4}
  \nonumber \\
  &= -\frac{9.87 \times 3.165 \times 10^{-26}}
  {720 \times 10^{-28}}
  \nonumber \\
  &= -\frac{3.12 \times 10^{-25}}{7.2 \times 10^{-26}}
  \nonumber \\
  &\approx -4.3\;\mathrm{J/m^3}
  = -4.8 \times 10^{-17}\;\mathrm{kg/m^3}\,.
  \label{eq:rho-casimir-num}
\end{align}
The gravitational $\DSigma$ induced over the plate separation is
\begin{equation}
  \DSigma_{\mathrm{Cas}} \sim \frac{8\pi G\,|\rhoCas|\,d^2}{c^4}\,.
  \label{eq:dsigma-casimir}
\end{equation}
Substituting:
\begin{align}
  \DSigma_{\mathrm{Cas}} &\sim
  \frac{8\pi \times 6.67 \times 10^{-11} \times 4.8 \times 10^{-17}
  \times 10^{-14}}{8.1 \times 10^{33}}
  \nonumber \\
  &\sim \frac{8.0 \times 10^{-40}}{8.1 \times 10^{33}}
  \nonumber \\
  &\sim 10^{-73}\,.
  \label{eq:dsigma-casimir-num}
\end{align}
This is extremely small.  Even using the Casimir energy over a
macroscopic stacking of $N = L/d$ plates over $L = 1\,\mathrm{cm}$
(so $N = 10^5$), the cumulative effect over the full stack is
\begin{equation}
  \DSigma_{\mathrm{stack}} \sim \frac{8\pi G\,|\rhoCas|\,L^2}{c^4}
  \sim \frac{8\pi \times 6.67 \times 10^{-11} \times 4.8 \times 10^{-17}
  \times 10^{-4}}{8.1 \times 10^{33}}
  \sim 10^{-63}\,.
  \label{eq:dsigma-stack}
\end{equation}

\textbf{Conclusion:} The Casimir energy density, while much larger than
the cosmological Khronon condensate (by $\sim 10$ orders of magnitude),
still produces a gravitational $\DSigma$ that is $\sim 48$ orders of
magnitude below the optical-clock threshold $10^{-15}$.  The Casimir
route is quantitatively excluded.

%======================================================================
\section{Approach 3: Superconductor Condensation}
\label{sec:superconductor}
%======================================================================

\subsection{Energy-based estimate}

A type-II superconductor undergoes a phase transition with condensation
energy density \known~\cite{Tinkham2004}
\begin{equation}
  u_{\mathrm{cond}} = \frac{B_c^2}{2\mu_0}\,,
  \label{eq:u-cond}
\end{equation}
where $B_c$ is the thermodynamic critical field.  For
$\mathrm{YBa_2Cu_3O_{7-\delta}}$ (YBCO) with
$B_c \sim 1\,\mathrm{T}$:
\begin{equation}
  u_{\mathrm{cond}} \sim \frac{1}{2 \times 4\pi \times 10^{-7}}
  \approx 4 \times 10^5\;\mathrm{J/m^3}\,.
  \label{eq:u-cond-num}
\end{equation}
The corresponding mass-energy density is
$\rho_{\mathrm{cond}} = u_{\mathrm{cond}}/c^2
\approx 4.4 \times 10^{-12}\;\mathrm{kg/m^3}$.

For a superconducting sample of size $L = 1\,\mathrm{cm}$:
\begin{align}
  \DSigma_{\mathrm{SC}} &\sim \frac{8\pi G\,\rho_{\mathrm{cond}}\,L^2}{c^4}
  \nonumber \\
  &= \frac{8\pi \times 6.67 \times 10^{-11} \times 4.4 \times 10^{-12}
  \times 10^{-4}}{8.1 \times 10^{33}}
  \nonumber \\
  &\sim \frac{7.4 \times 10^{-25}}{8.1 \times 10^{33}}
  \sim 10^{-58}\,.
  \label{eq:dsigma-sc-energy}
\end{align}

\noindent More precisely, the \emph{change} in stress-energy upon
entering the superconducting state is $\Delta T_{00} = u_{\mathrm{cond}}$,
and the metric perturbation it sources is of order
$h_{00} \sim 8\pi G\,u_{\mathrm{cond}}\,L^2/c^4$, yielding the same
estimate \derived.  This is $\sim 43$ orders below threshold.

\subsection{Quantum-coherence hypothesis}
\label{subsec:coherence}

A more speculative avenue exists \specmark.  In the normal metallic
state, conduction electrons undergo frequent scattering events, each
producing irreversibility (entropy production).  In the
$\Sigma$--$\tauasym$ framework, each scattering event contributes
$\delta\Sigma > 0$.  In the superconducting state, Cooper pairs
propagate without scattering: $\tauasym_{\mathrm{electron}} = 0$.

If the \emph{gravitational} $\Sgrav$ couples to the
\emph{microscopic} entropy production of matter (beyond the
stress-energy tensor alone), then the normal-to-superconductor
transition eliminates a source term for $\Sgrav$.

We estimate the electron scattering contribution to
local entropy production.  For copper at $T = 300\,\mathrm{K}$
\known:
\begin{itemize}
  \item Electron density: $n_e \approx 8.5 \times 10^{28}\,\mathrm{m^{-3}}$
  \item Fermi velocity: $\vF \approx 1.6 \times 10^6\,\mathrm{m/s}$
  \item Relaxation time: $\trel \approx 2.5 \times 10^{-14}\,\mathrm{s}$
  \item Mean free path: $\ell_{\mathrm{mfp}} = \vF\,\trel
    \approx 4 \times 10^{-8}\,\mathrm{m}$
\end{itemize}
The entropy produced per electron per scattering is of order
\begin{equation}
  \delta\Sigma_{\mathrm{scatter}} \sim \frac{G\,m_e^2}{\hbar\,c\,
  \ell_{\mathrm{mfp}}}
  \label{eq:sigma-scatter}
\end{equation}
where $m_e = 9.11 \times 10^{-31}\,\mathrm{kg}$.  This is the
gravitational self-energy of the electron over one mean free path,
in natural units:
\begin{align}
  \delta\Sigma_{\mathrm{scatter}} &\sim
  \frac{6.67 \times 10^{-11} \times (9.11 \times 10^{-31})^2}
  {1.055 \times 10^{-34} \times 3 \times 10^8 \times 4 \times 10^{-8}}
  \nonumber \\
  &= \frac{6.67 \times 10^{-11} \times 8.3 \times 10^{-61}}
  {1.27 \times 10^{-33}}
  \nonumber \\
  &= \frac{5.5 \times 10^{-71}}{1.27 \times 10^{-33}}
  \approx 4.4 \times 10^{-38}\,.
  \label{eq:sigma-scatter-num}
\end{align}
The scattering rate per electron is $1/\trel \approx 4 \times 10^{13}\,
\mathrm{s^{-1}}$.  In a sample volume $V = (1\,\mathrm{cm})^3 = 10^{-6}\,
\mathrm{m^3}$, the total entropy production rate is
\begin{align}
  \dot{\Sigma}_{\mathrm{total}} &= n_e \times V \times
  \frac{\delta\Sigma_{\mathrm{scatter}}}{\trel}
  \nonumber \\
  &= 8.5 \times 10^{28} \times 10^{-6} \times
  \frac{4.4 \times 10^{-38}}{2.5 \times 10^{-14}}
  \nonumber \\
  &= 8.5 \times 10^{22} \times 1.76 \times 10^{-24}
  \nonumber \\
  &\approx 0.15\;\mathrm{s^{-1}}\,.
  \label{eq:sigma-dot}
\end{align}

However, this is the gravitational entropy production
\emph{rate} from electron-scattering self-gravity, not the
\emph{steady-state} $\DSigma$ sourced by the scattering.  The
steady-state requires solving the Poisson-like equation
$\nabla^2(\DSigma) = S$ with appropriate boundary conditions.
The source $S \sim G\,n_e\,m_e^2/(\hbar\,c\,\trel)$
produces a static $\DSigma$ of order
\begin{equation}
  \DSigma_{\mathrm{coherence}} \sim S \times L^2
  \sim \frac{G\,n_e\,m_e^2\,L^2}{\hbar\,c\,\trel}\,.
  \label{eq:dsigma-coherence}
\end{equation}
Substituting:
\begin{align}
  \DSigma_{\mathrm{coherence}} &\sim
  \frac{6.67 \times 10^{-11} \times 8.5 \times 10^{28}
  \times 8.3 \times 10^{-61} \times 10^{-4}}
  {1.055 \times 10^{-34} \times 3 \times 10^8
  \times 2.5 \times 10^{-14}}
  \nonumber \\
  &= \frac{6.67 \times 8.5 \times 8.3 \times 10^{-148}}
  {7.9 \times 10^{-41}}
  \nonumber \\
  &\sim \frac{4.7 \times 10^{-146}}{7.9 \times 10^{-41}}
  \sim 6 \times 10^{-106}\,.
  \label{eq:dsigma-coherence-num}
\end{align}

This is \emph{far} below any detection threshold, essentially because
the electron gravitational self-energy is negligibly small compared
to its electromagnetic interactions.  The often-cited claims of
``gravitational effects in superconductors''
\cite{Podkletnov1992,Tajmar2006} would require $\DSigma \sim 10^{-8}$
to $10^{-5}$, which is $\sim 100$ orders of magnitude above what the
electron self-gravity channel can produce.

\textbf{Conclusion:} The energy-based estimate gives $\DSigma \sim
10^{-58}$; the quantum-coherence route gives $\DSigma \sim 10^{-106}$.
Neither is detectable.  However, we stress that
Eq.~\eqref{eq:dsigma-coherence} assumes $\Sgrav$ couples to matter
only through gravitational self-energy.  If a \emph{non-perturbative}
coupling between macroscopic quantum coherence and $\Sgrav$ exists
(for instance, via the topological properties of the Cooper-pair
condensate), the estimate could change qualitatively.  This remains
an open question.

%======================================================================
\section{Approach 4: Nanoparticle Interferometry}
\label{sec:nanoparticle}
%======================================================================

Recent experiments have achieved matter-wave interference with
increasingly massive particles, culminating in the demonstration
by Pedalino \textit{et al.}~\cite{Pedalino2026} of interference
fringes for nanoparticles of mass $m = 170\,\mathrm{kDa}
= 2.8 \times 10^{-22}\,\mathrm{kg}$.

When a massive particle is placed in a spatial superposition
\begin{equation}
  \ket{\psi} = \frac{1}{\sqrt{2}}
  \bigl(\ket{\bm{r}_1} + \ket{\bm{r}_2}\bigr)
  \label{eq:superposition}
\end{equation}
with separation $\Delta x = |\bm{r}_1 - \bm{r}_2|$, the
gravitational self-interaction produces an entropy production
\begin{equation}
  \Sigma_{\mathrm{self}} = \frac{G\,m^2}{\hbar\,c\,\Delta x}\,.
  \label{eq:sigma-self}
\end{equation}
This can be understood as follows \derived: the gravitational potential
energy between the two branches is $E_{\mathrm{grav}} = G m^2/\Delta x$,
and the associated $\Sgrav$ is $E_{\mathrm{grav}}/E_{\mathrm{Planck}}$
where $E_{\mathrm{Planck}} = \hbar c / L_{\mathrm{Planck}}^2 \times L_{\mathrm{Planck}}
= \sqrt{\hbar c^5/G}$.  More precisely:
\begin{equation}
  \Sigma_{\mathrm{self}} = \frac{G\,m^2}{\hbar\,c\,\Delta x}
  = \frac{t_{\mathrm{grav}}}{t_{\mathrm{Planck}}} \times
  \frac{L_{\mathrm{Planck}}}{\Delta x}\,,
  \label{eq:sigma-self-interpret}
\end{equation}
where $t_{\mathrm{grav}} = Gm/(c^2 \Delta x) \times \Delta x/c$ is
the gravitational light-crossing time.

\subsection{Current experiments}

For the Pedalino \textit{et al.} parameters
($m = 2.8 \times 10^{-22}\,\mathrm{kg}$,
$\Delta x \sim 10\,\mathrm{nm} = 10^{-8}\,\mathrm{m}$):
\begin{align}
  \Sigma_{\mathrm{self}} &=
  \frac{6.67 \times 10^{-11} \times (2.8 \times 10^{-22})^2}
  {1.055 \times 10^{-34} \times 3 \times 10^8 \times 10^{-8}}
  \nonumber \\
  &= \frac{6.67 \times 10^{-11} \times 7.84 \times 10^{-44}}
  {3.165 \times 10^{-34}}
  \nonumber \\
  &= \frac{5.23 \times 10^{-54}}{3.165 \times 10^{-34}}
  \nonumber \\
  &\approx 1.7 \times 10^{-20}\,.
  \label{eq:sigma-current}
\end{align}
This is too small to cause observable decoherence on experimental
timescales (the decoherence rate scales as
$\Gamma \sim \Sigma_{\mathrm{self}} \times c/\Delta x$
\cite{Diosi1989,Penrose1996}).

\subsection{Next-generation experiments}

Proposals for next-generation experiments target masses
$m \sim 10^9\,\mathrm{amu} = 1.7 \times 10^{-18}\,\mathrm{kg}$
and superposition separations $\Delta x \sim 1\,\mu\mathrm{m}
= 10^{-6}\,\mathrm{m}$ \cite{Bose2017,Marletto2017}:
\begin{align}
  \Sigma_{\mathrm{self}}^{\mathrm{next}} &=
  \frac{6.67 \times 10^{-11} \times (1.7 \times 10^{-18})^2}
  {1.055 \times 10^{-34} \times 3 \times 10^8 \times 10^{-6}}
  \nonumber \\
  &= \frac{6.67 \times 10^{-11} \times 2.89 \times 10^{-36}}
  {3.165 \times 10^{-32}}
  \nonumber \\
  &= \frac{1.93 \times 10^{-46}}{3.165 \times 10^{-32}}
  \nonumber \\
  &\approx 6.1 \times 10^{-15}\,.
  \label{eq:sigma-next}
\end{align}
This is tantalizingly close to the optical-clock threshold
$\DSigma \sim 10^{-15}$ and well within the range that could
produce \emph{measurable gravitational decoherence}.

The Bose--Marletto--Vedral (BMV) proposal \cite{Bose2017,Marletto2017}
provides a direct experimental test: two particles, each in spatial
superposition, interact gravitationally.  If they become entangled,
this certifies that gravity mediates quantum information---and in our
framework, that $\Sigma_{\mathrm{self}}$ is a quantum observable.
The entanglement phase is
\begin{equation}
  \phi_{\mathrm{ent}} = \frac{G\,m^2\,t}{\hbar\,\Delta x}
  = \Sigma_{\mathrm{self}} \times \frac{c\,t}{\Delta x}\,.
  \label{eq:bose-marletto}
\end{equation}
For $\phi_{\mathrm{ent}} \sim 1$ (detectable), one needs
$t \sim \Delta x/(c\,\Sigma_{\mathrm{self}})$.  With the
next-generation parameters:
\begin{equation}
  t \sim \frac{10^{-6}}{3 \times 10^8 \times 6 \times 10^{-15}}
  \sim 0.6\;\mathrm{s}\,,
  \label{eq:bose-marletto-time}
\end{equation}
which is within the reach of proposed experiments
\cite{Bose2017,VanDeKamp2020}.

\subsection{Interpretation within the $\Sigma$ framework}

If the BMV experiment succeeds, it demonstrates that $\Sgrav$ at the
scale $\Sigma \sim 10^{-14}$ is a genuine quantum-information
quantity, not merely a classical entropy.  This would validate the
central hypothesis of our framework---that $\Sgrav = -\ln(-g_{00})$
is a fundamental field obeying quantum mechanics---at a scale where
its effects are directly measurable.

Conversely, if gravitational decoherence is observed at a rate
\emph{exceeding} $\Gamma = G m^2/(2\hbar\,\Delta x^2)$
(the Di\'osi--Penrose prediction), this would indicate additional
non-unitary contributions to $\Sgrav$ beyond the self-gravitational
channel, potentially pointing to new physics.

%======================================================================
\section{Feasibility Summary}
\label{sec:summary}
%======================================================================

Table~\ref{tab:summary} collects our estimates.

\begin{table*}[t]
\caption{Feasibility of local $\Sigma$ modification by four approaches.
All numerical values are order-of-magnitude estimates.  The detection
thresholds assume current best technology; ``next-gen'' refers to
experiments planned for the 2030s.
\label{tab:summary}}
\begin{ruledtabular}
\begin{tabular}{lcccl}
Approach & $\DSigma$ achievable & Best threshold & Gap (orders)
& Status \\
\colrule
Khronon condensate (\S\ref{sec:khronon})
  & $10^{-52}$ & $10^{-15}$ & 37 & Excluded \\
Casimir vacuum (\S\ref{sec:casimir})
  & $10^{-63}$ & $10^{-15}$ & 48 & Excluded \\
Superconductor, energy (\S\ref{sec:superconductor})
  & $10^{-58}$ & $10^{-15}$ & 43 & Excluded \\
Superconductor, coherence (\S\ref{subsec:coherence})
  & $10^{-106}$ & $10^{-15}$ & 91 & Excluded (but see text) \\
Nanoparticle, current (\S\ref{sec:nanoparticle})
  & $10^{-20}$ & $10^{-20}$ & 0 & Marginal \\
Nanoparticle, next-gen (\S\ref{sec:nanoparticle})
  & $10^{-14}$ & $10^{-15}$ & $-1$ (above) & \textbf{Feasible} \\
\end{tabular}
\end{ruledtabular}
\end{table*}

%======================================================================
\section{Discussion}
\label{sec:discussion}
%======================================================================

Our analysis leads to three principal conclusions:

\textit{1. Laboratory gravity modification is excluded by many orders
of magnitude.}---None of the approaches we considered can produce a
detectable $\DSigma$ at a location devoid of additional mass.  The
fundamental obstacle is the weakness of the gravitational coupling:
$G/c^4 \approx 8.3 \times 10^{-45}\,\mathrm{m/J}$, which converts
even large energy densities into minuscule metric perturbations.
Claims of anomalous gravitational effects in superconductors
\cite{Podkletnov1992,Tajmar2006} would require $\DSigma$ values
exceeding our most optimistic estimates by $40$--$100$ orders of
magnitude and should be regarded as unsubstantiated.

\textit{2. Nanoparticle interferometry is the most promising path to
detecting $\Sgrav$ at the quantum scale.}---Next-generation
experiments with $m \sim 10^9\,\mathrm{amu}$ and
$\Delta x \sim 1\,\mu\mathrm{m}$ achieve
$\Sigma_{\mathrm{self}} \sim 10^{-14}$, which is \emph{above} the
threshold for observable gravitational decoherence and entanglement.
The BMV experiment \cite{Bose2017,Marletto2017} would directly test
whether $\Sgrav$ is a quantum observable, providing the first
empirical anchor for the $\Sigma$--$\tauasym$ framework at the
quantum-gravitational interface.

\textit{3. The superconductor-coherence channel deserves further
theoretical investigation.}---Our estimate in
\S\ref{subsec:coherence} assumed that $\Sgrav$ couples to matter
only through the gravitational self-energy of individual electrons.
If macroscopic quantum coherence modifies the \emph{boundary
conditions} of $\Sgrav$ (rather than sourcing it directly), the
relevant quantity might be the topological winding number of the
order parameter, not the condensation energy.  This possibility
cannot be excluded by the present analysis and merits a dedicated
study within the framework of Paper~6 \cite{Huang2026paper6},
where the gauge structure of the $\Sigma$ field is developed.

Finally, we note an important conceptual point.  Even though
``gravity control'' in the engineering sense (generating macroscopic
$\DSigma$ without mass) appears infeasible, the $\Sigma$ framework
provides a new \emph{diagnostic language} for gravitational physics:
every gravitational measurement is a measurement of $\Sgrav$, and
the Petz bound~\eqref{eq:petz-bound} provides a universal lower
limit on the achievable recovery fidelity.  This
information-theoretic perspective may prove as valuable for
gravitational-wave astronomy and precision geodesy as it is for
the foundational questions addressed in this paper.

%======================================================================
\begin{acknowledgments}
The author thanks the developers of the Petz recovery map formalism
and the experimental groups pushing the mass frontier of matter-wave
interferometry.
\end{acknowledgments}
%======================================================================

\begin{thebibliography}{30}

\bibitem{Huang2026paper1}
S.-K.~Huang,
``Temporal asymmetry as Petz recovery failure: A unified
information-theoretic arrow of time,''
(2026), Paper~1 of this series.

\bibitem{Petz1986}
D.~Petz,
``Sufficient subalgebras and the relative entropy of states of a
von~Neumann algebra,''
Commun.\ Math.\ Phys.\ \textbf{105}, 123 (1986).

\bibitem{Petz1988}
D.~Petz,
``Sufficiency of channels over von~Neumann algebras,''
Q.\ J.\ Math.\ \textbf{39}, 97 (1988).

\bibitem{Junge2018}
M.~Junge, R.~Renner, D.~Sutter, M.~M.~Wilde, and A.~Winter,
``Universal recovery maps and approximate sufficiency of quantum
relative entropy,''
Ann.\ Henri Poincar\'e \textbf{19}, 2955 (2018).

\bibitem{Faulkner2022}
T.~Faulkner, S.~Hollands, B.~Swingle, and Y.~Wang,
``Approximate recovery and relative entropy I: general von~Neumann
subalgebras,''
Commun.\ Math.\ Phys.\ \textbf{389}, 349 (2022).

\bibitem{Huang2026paper2}
S.-K.~Huang,
``The gravitational refractive index as Petz recovery fidelity:
Temporal asymmetry from spacetime geometry,''
(2026), Paper~2 of this series.

\bibitem{Adelberger2009}
E.~G.~Adelberger, J.~H.~Gundlach, B.~R.~Heckel, S.~Hoedl, and
S.~Schlamminger,
``Torsion balance experiments: A low-energy frontier of particle
physics,''
Prog.\ Part.\ Nucl.\ Phys.\ \textbf{62}, 102 (2009).

\bibitem{Bothwell2022}
T.~Bothwell \textit{et al.},
``Resolving the gravitational redshift across a millimetre-scale
atomic sample,''
Nature \textbf{602}, 420 (2022).

\bibitem{Asenbaum2020}
P.~Asenbaum, C.~Overstreet, M.~Kim, J.~Curti, and M.~A.~Kasevich,
``Atom-interferometric test of the equivalence principle at the
$10^{-12}$ level,''
Phys.\ Rev.\ Lett.\ \textbf{125}, 191101 (2020).

\bibitem{LIGOScientific2015}
B.~P.~Abbott \textit{et al.} (LIGO Scientific and Virgo Collaborations),
``Observation of gravitational waves from a binary black hole merger,''
Phys.\ Rev.\ Lett.\ \textbf{116}, 061102 (2016).

\bibitem{Huang2026paper3}
S.-K.~Huang,
``Dark matter as Khronon condensate: Weak-field phenomenology of the
$\tau$ framework,''
(2026), Paper~3 of this series.

\bibitem{ArkaniHamed2004}
N.~Arkani-Hamed, H.-C.~Cheng, M.~A.~Luty, and S.~Mukohyama,
``Ghost condensation and a consistent infrared modification of
gravity,''
J.\ High Energy Phys.\ \textbf{05} (2004) 074.

\bibitem{Casimir1948}
H.~B.~G.~Casimir,
``On the attraction between two perfectly conducting plates,''
Proc.\ K.\ Ned.\ Akad.\ Wet.\ \textbf{51}, 793 (1948).

\bibitem{Lamoreaux1997}
S.~K.~Lamoreaux,
``Demonstration of the Casimir force in the 0.6 to $6\,\mu$m range,''
Phys.\ Rev.\ Lett.\ \textbf{78}, 5 (1997).

\bibitem{Tinkham2004}
M.~Tinkham,
\textit{Introduction to Superconductivity}, 2nd ed.\
(Dover, New York, 2004).

\bibitem{Podkletnov1992}
E.~Podkletnov and R.~Nieminen,
``A possibility of gravitational force shielding by bulk
$\mathrm{YBa_2Cu_3O_{7-x}}$ superconductor,''
Physica~C \textbf{203}, 441 (1992).

\bibitem{Tajmar2006}
M.~Tajmar, F.~Plesescu, B.~Seifert, R.~Schnitzer, and I.~Vasiljevich,
``Search for frame-dragging-like signals close to spinning
superconductors,''
J.\ Phys.:\ Conf.\ Ser.\ \textbf{150}, 032101 (2009).

\bibitem{Pedalino2026}
S.~Pedalino \textit{et al.},
``Matter-wave interference of 170-kDa nanoparticles,''
Nature \textbf{649}, 866 (2026).

\bibitem{Diosi1989}
L.~Di\'osi,
``Models for universal reduction of macroscopic quantum
fluctuations,''
Phys.\ Rev.\ A \textbf{40}, 1165 (1989).

\bibitem{Penrose1996}
R.~Penrose,
``On gravity's role in quantum state reduction,''
Gen.\ Relativ.\ Gravit.\ \textbf{28}, 581 (1996).

\bibitem{Bose2017}
S.~Bose \textit{et al.},
``Spin entanglement witness for quantum gravity,''
Phys.\ Rev.\ Lett.\ \textbf{119}, 240401 (2017).

\bibitem{Marletto2017}
C.~Marletto and V.~Vedral,
``Gravitationally induced entanglement between two massive particles
is sufficient evidence of quantum effects of gravity,''
Phys.\ Rev.\ Lett.\ \textbf{119}, 240402 (2017).

\bibitem{VanDeKamp2020}
R.~J.~van de Kamp \textit{et al.},
``Quantum gravity witness via entanglement of masses: Casimir
screening,''
Phys.\ Rev.\ A \textbf{102}, 062807 (2020).

\bibitem{Huang2026paper6}
S.-K.~Huang,
``Electromagnetic force from temporal asymmetry: U(1) gauge structure
and DBI-$\Sigma$ correspondence,''
(2026), Paper~6 of this series.

\end{thebibliography}

\end{document}
